% !TeX root = ../数模通用模板.tex
\begin{abstract}
分布式光伏为社区提供了就地获取电能的途径，但其电力供给的时间分布与居民用电需求并不天然一致，负载、光伏与电价的变化使微网购电同时面临费用控制、供电保障和储能运行强度的权衡。本文建立了一个“预测—计划—修正—执行—结算”一体化的调度分析模型，对建设兼具经济性与可靠性的微网发电系统具有重要意义。

\textbf{对于问题一，}为了体现现实中储能电池不能碎片化、剧烈改变其充放电功率，建立两层混合整数规划，在功率爬坡、最短开启与关闭时间等约束下，先求最低购电费，再允许费用增加0.1\%以减小功率总变差。典型日第二层购电量为57560.6438\,kWh，费用为34115.03元，初末储电量均为6000\,kWh。相对只考虑经济性的基准，费用增加0.93\%，功率总变差下降67.66\%，充放电开启次数由12次减至4次。

\textbf{对于问题二，}先等权融合直方图梯度提升树与极端随机树，再进行共同岭回归校准和前日负载偏差半量修正。利用最近28日的整日预测误差保留时间关联，依次执行模式规划、固定模式购电修正和当槽储能反馈。评价期总购电费为13195653.64元，实际非空充放电方向反转2535次；每日计划模式变化预算与实际反转次数分别核验，不将二者混同。

\textbf{对于问题三，}对官方光伏预报作同发布时间节点校准，以负载已观察前缀修正日内偏差，构造共享购电量的情景线性规划。模型在0、6、12、18时更新剩余计划，下调按取消部分退回原款一半，并按相对午夜计划的最终净差额结算。评价期费用为12788960.26元，其中下调退款77458.97元，紧急购电费202487.48元。

\textbf{对于问题四，}沿用问题二的预测思路，按价格随时间揭示的条件分别滚动预测未来电价，将无日内更新和有日内更新的完整方案用于同一波动价格序列。两方案评价期费用分别为14008609.54元和13539823.62元，后者低3.35\%。

最终，我们使用敏感性分析和固定计划扰动回放等方式验证模型的泛化能力，并进一步说明结果的适用范围。

\keywords{微网购电；储能调度；混合整数规划；情景线性规划；预测校准}
\end{abstract}
